% 3.5 章 · 动力系统、分岔与混沌
% 这一章的任务只有一个：把"混沌"这三个字换成**可以量出来的数**。
%   1) 指数分离率（Lyapunov 指数）——配"重正交化 QR + 变分方程"的口径与步长证据；
%   2) 分岔结构（周期倍增、Feigenbaum 常数、周期窗口）——配参数网格分辨率；
%   3) 几何（吸引子、Poincare 截面、零倾线）——配"轨道数、丢弃的暂态、投影参数"。
% 三条都要给出**实测数字**，并且都要能回答"这个数字凭什么不是数值噪声"。
% 断言附表由 tools/build_book.py 从实际运行结果生成，\input 进来，不手抄。

\numlabchapter{动力系统、分岔与混沌}{3.5}

\section{这一章要解决什么}

"混沌"这个词的危险之处在于它听起来像一个形容词。本章把它拆成三件可测量的事：

\begin{enumerate}
  \item \textbf{指数分离率}。两条初值相差 $\varepsilon$ 的轨道，平均而言按
        $d(t)\approx\varepsilon e^{\lambda_1 t}$ 分开。$\lambda_1>0$ 就是"对初值敏感"，
        它是本章第一个可以直接报数的量。
  \item \textbf{分岔结构}。参数连续变化时，周期轨在什么参数值上出现、失稳、倍增、
        消失。这些都是方程（乘子 $=+1$ 或 $-1$）的根，可以算到机器精度。
  \item \textbf{几何}。吸引子长什么样、轨道多久回来一次、零倾线把平面切成几块。
        这一部分只能用图来说，但图上的每个点仍然必须来自解出来的方程，
        而不是"颜色看着像"。
\end{enumerate}

贯穿全章的一条纪律：\textbf{任何"混沌"结论都必须附一条数值证据链}。
本章用的是三条腿：散度恒等式（$\sum\lambda_i=\langle\nabla\!\cdot\!f\rangle$）、
步长减半（$h\to h/2$ 后结论不变）、以及"解析对拍"（能闭式的一律对拍）。
记号：$u=\tfrac12\varepsilon_{\text{mach}}=2^{-53}\approx1.11\times10^{-16}$；
本章所有"实测"数字都来自 \texttt{tests/3.5/} 的断言报告，
"本仓实测"指的是固定口径下的复现结果，口径随数字一起给出。

\section{Lyapunov 指数：把"敏感"变成一个数}

\begin{definition}[流的 Lyapunov 指数]
设 $f$ 生成流 $\varphi_t$，$D\varphi_t(y_0)$ 是沿轨道 $y(t)=\varphi_t(y_0)$ 的变分（切）映射，
即 $D\varphi_t$ 满足
\begin{equation}
  \frac{d}{dt}D\varphi_t=J(y(t))\,D\varphi_t,\qquad D\varphi_0=I ,
  \label{eq:variational}
\end{equation}
其中 $J=\partial f/\partial y$。把 $D\varphi_t$ 的奇异值记为 $\sigma_1(t)\ge\sigma_2(t)\ge\dots$，
则第 $i$ 个 Lyapunov 指数定义为
\begin{equation}
  \lambda_i=\lim_{t\to\infty}\frac1t\log\sigma_i(t).
  \label{eq:lyap}
\end{equation}
$\lambda_1>0$ 时称轨道（在该不变集上）是混沌的。
\end{definition}

对映射 $x_{n+1}=g(x_n)$，把 $D\varphi_t$ 换成 $\prod_{k=0}^{n-1}Dg(x_k)$ 即可，
下文所有结论对映射逐字成立。

\begin{theorem}[指数之和等于散度的时间平均]
\label{thm:sum-div}
对 $\mathbb{R}^d$ 上由 $f$ 生成的流，
\begin{equation}
  \sum_{i=1}^{d}\lambda_i=\lim_{T\to\infty}\frac1T\int_0^T\nabla\!\cdot\!f(y(t))\,dt .
  \label{eq:sumdiv}
\end{equation}
特别地，若 $\nabla\!\cdot\!f$ 是常数 $c$，则 $\sum_i\lambda_i=c$。
对映射，$\sum_i\lambda_i=\lim_{n\to\infty}\frac1n\sum_{k<n}\log|\det Dg(x_k)|$。
\end{theorem}

\begin{proof}
设 $M(t)=D\varphi_t$。由 \eqref{eq:variational} 与 Jacobi 公式
$\frac{d}{dt}\det M=(\operatorname{tr}J)\det M$（在 $\det M>0$ 的邻域上对
$\log\det M$ 求导即可），得
\[
  \det D\varphi_t=\exp\!\Bigl(\int_0^t\nabla\!\cdot\!f(y(s))\,ds\Bigr).
\]
另一方面 $\det D\varphi_t=\prod_i\sigma_i(t)$，于是
\[
  \frac1t\log\bigl|\det D\varphi_t\bigr|=\sum_i\frac1t\log\sigma_i(t)
  =\frac1t\int_0^t\nabla\!\cdot\!f(y(s))\,ds .
\]
两边令 $t\to\infty$，左边逐项收敛到 $\sum_i\lambda_i$，右边收敛到散度的时间平均。
常数散度的情形立刻得到。映射的情形把行列式换成乘积、把积分换成求和即可
（$\log|\det\prod Dg|=\sum_k\log|\det Dg(x_k)|$）。
\end{proof}

这条定理是本章的\textbf{第一根支柱}，因为它给出一个不依赖任何文献值的标尺：

\begin{itemize}
  \item Lorenz 的散度处处是 $-\left(\sigma+1+\beta\right)$。取 $\sigma=10$、$\beta=8/3$ 得
        $-13.6666666667$。本仓实测（$T=2000$、$h=0.01$、QR 每 $10$ 步、暂态 $50$）
        两条初值的 $\sum\lambda$ 与它的差都小于 $5\times10^{-4}$。
  \item Henon 映射的 $|\det Dg|=|b|$ 处处常数，$b=0.3$ 时
        $\sum\lambda=\log 0.3=-1.2039728043259361$。本仓实测
        $\lambda=(0.4193,-1.6233)$，和与闭式的差小于 $10^{-12}$。
\end{itemize}

\begin{theorem}[Benettin 的重正交化 QR 给出全部指数]
\label{thm:benettin}
设 $Q_nR_n=D\varphi_{n\Delta t}$ 是每 $\Delta t$ 步做一次 QR 分解的结果，
$R_n$ 取正对角。则
\begin{equation}
  \lambda_i=\lim_{n\to\infty}\frac{1}{n\Delta t}\sum_{k=1}^{n}\log (R_k)_{ii},
  \label{eq:benettin}
\end{equation}
且该极限与每次 QR 的间隔 $\Delta t$、与正交化的具体算法（Gram--Schmidt 或 Householder）
无关。
\end{theorem}

\begin{proof}[证明要点]
把 $D\varphi_{n\Delta t}$ 写成相邻块的乘积
$M_nM_{n-1}\cdots M_1$（$M_k=D\varphi_{k\Delta t\to(k+1)\Delta t}$），
并归纳地记 $M_kQ_{k-1}=Q_kR_k$。这就是 \eqref{eq:benettin} 的递推：
每一步先传播、再正交化。逐块 QR 与整体 QR 的差别只在于 $R$ 的对角元，
而对角元满足 $(R_k)_{ii}=\|(I-P_{i-1})M_kq^{(k-1)}_i\|$，
其中 $P_{i-1}$ 是前 $i-1$ 个方向张成的正交投影——
这正是"第 $i$ 个方向在每一步里的伸长倍数"。
Gram--Schmidt 与 Householder 给出同一组 $\{|(R_k)_{ii}|\}$（差一个符号，
由正对角约定消掉），所以指数与算法无关。
会改变结果的是 $\Delta t$：只有当每一步的伸长倍数不溢出、
且重正交化的次数足够多时，递推才等价于奇异值的极限；
数值上表现为\textbf{正交性缺陷}必须留在舍入水平。
本仓实测（Lorenz，$T=2000$，$300$ 次重正交化）最大缺陷 $4.4\times10^{-16}$；
把间隔从 $10$ 步改成 $20$ 步，谱的变化小于 $10^{-9}$——
这条"间隔无关性"就是定理里那句"与 $\Delta t$ 无关"的实测版本。
\end{proof}

\begin{corollary}[Kaplan--Yorke 维数]
设 $\lambda_1\ge\dots\ge\lambda_d$，取最大的 $j$ 使 $\sum_{i\le j}\lambda_i\ge0$，则
\begin{equation}
  D_{KY}=j+\frac{\sum_{i\le j}\lambda_i}{|\lambda_{j+1}|}.
  \label{eq:ky}
\end{equation}
Lorenz 用本仓实测谱代入得 $D_{KY}=2.062$，与文献常用的 $2.06$ 一致。
注意这一条\textbf{不是独立证据}：它由同一组指数推出，只能作一致性检查。
\end{corollary}

\section{Lorenz：三个能闭式对账的量}

\begin{theorem}[平衡点、散度与相型]
\label{thm:lorenz-fp}
对 $\dot x=\sigma(y-x)$、$\dot y=x(\rho-z)-y$、$\dot z=xy-\beta z$：
\begin{enumerate}
  \item 三个平衡点是 $O=(0,0,0)$ 与
        $C_\pm=\bigl(\pm\sqrt{\beta(\rho-1)},\pm\sqrt{\beta(\rho-1)},\rho-1\bigr)$
        （$\rho>1$ 时存在）；
  \item $\nabla\!\cdot\!f=-(\sigma+1+\beta)$ 与状态无关；
  \item $O$ 的特征值是 $\lambda^3+(\sigma+\beta+1)\lambda^2+(\beta+\rho\sigma)\lambda
        +\sigma\beta(1-\rho)$ 的根；$C_\pm$ 的特征值满足
        \begin{equation}
          \lambda^3+(\sigma+\beta+1)\lambda^2+\beta(\rho+\sigma)\lambda
          +2\sigma\beta(\rho-1)=0 .
          \label{eq:lorenz-char}
        \end{equation}
\end{enumerate}
\end{theorem}

\begin{proof}
(1) 令右端为零：$\dot x=0$ 给 $y=x$；$\dot z=0$ 给 $xy=\beta z$，故 $z=x^2/\beta$；
代入 $\dot y=0$ 得 $x(\rho-z)-x=0$，即 $x\bigl(\rho-1-x^2/\beta\bigr)=0$，
于是 $x=0$（$O$）或 $x^2=\beta(\rho-1)$（$C_\pm$）。
(2) 直接求迹：$-\sigma-1-\beta$。
(3) 在 $C_\pm$ 处
\[
  J=\begin{pmatrix}-\sigma&\sigma&0\\ 1&-1&-x\\ x&x&-\beta\end{pmatrix},
  \qquad x=\sqrt{\beta(\rho-1)},
\]
展开 $\det(\lambda I-J)$ 并用 $x^2=\beta(\rho-1)$ 化简即得 \eqref{eq:lorenz-char}；
在 $O$ 处 $x=0$，同一展开给 (3) 中写出的三次式。
\end{proof}

$\rho=28$（经典参数）时：$O$ 的特征值是 $(-22.8277,-2.6667,11.8277)$，
一个正实根加两个负实根；$C_\pm$ 的是 $(-13.8546,\,0.0940\pm10.1945\mathrm{i})$。
注意后者：\textbf{一个稳定实方向 $+$ 一个不稳定螺旋面}。
按相型分类这是 \texttt{saddle\_focus}（鞍焦），不是 \texttt{saddle}（鞍点）——
鞍点的两个方向都是实的。本仓初版把所有"有正有负"的特征值都判成鞍点，
于是 $\rho=28$ 的 $C_\pm$ 被写成了鞍点；这条分类现在是断言 \texttt{test\_lorenz\_}\allowbreak
\texttt{fixed\_}\allowbreak\texttt{points\_}\allowbreak\texttt{and\_}\allowbreak
\texttt{divergence} 的回归项。这个区别不是术语洁癖：奇异吸引子的两片"翅膀"
正是绕着 $C_+$ 与 $C_-$ 的两个螺旋面转出来的。

\begin{theorem}[Hopf 阈值 $\rho_H=470/19$]
\label{thm:lorenz-hopf}
在 \eqref{eq:lorenz-char} 中令 $\lambda=\mathrm{i}\omega$（$\omega>0$），
方程组有解当且仅当
\begin{equation}
  \rho=\rho_H=\frac{\sigma(\sigma+\beta+3)}{\sigma-\beta-1},\qquad
  \omega^2=\beta(\rho_H+\sigma)=\frac{2\sigma\beta(\rho_H-1)}{\sigma+\beta+1}.
  \label{eq:hopf}
\end{equation}
取 $\sigma=10,\beta=8/3$ 得 $\rho_H=470/19=24.7368421053$，$\omega^2=92.6315789474$。
\end{theorem}

\begin{proof}
把 $\lambda=\mathrm{i}\omega$ 代入 \eqref{eq:lorenz-char} 并分开实部、虚部：
\[
  \text{实部：}\ -(\sigma+\beta+1)\omega^2+2\sigma\beta(\rho-1)=0,
  \qquad
  \text{虚部：}\ \mathrm{i}\omega\bigl(\beta(\rho+\sigma)-\omega^2\bigr)=0 .
\]
由虚部（$\omega\ne0$）得 $\omega^2=\beta(\rho+\sigma)$；代入实部得
\[
  (\sigma+\beta+1)\beta(\rho+\sigma)=2\sigma\beta(\rho-1)
  \ \Longrightarrow\
  (\sigma+\beta+1)(\rho+\sigma)=2\sigma(\rho-1).
\]
展开并按 $\rho$ 归并：
$\rho\bigl(\sigma+\beta+1-2\sigma\bigr)=-\sigma(\sigma+\beta+1)-2\sigma$，
即 $\rho(\beta+1-\sigma)=-\sigma(\sigma+\beta+3)$，
两边同除 $\beta+1-\sigma=-(\sigma-\beta-1)$ 得 \eqref{eq:hopf}。
\end{proof}

本仓实测：用"$\max\operatorname{Re}\lambda(\rho)$ 变号"做判据二分 $40$ 次，
得 $\rho_H=24.736842105253$，与 $470/19$ 差 $9.8\times10^{-12}$。
在 $\rho_H$ 处 $C_\pm$ 的特征值是 $(-13.6667,\,0\pm9.6245\mathrm{i})$：
\textbf{第三个特征值不是零}（它等于 $-\left(\sigma+1+\beta\right)$）。
把判据写成 $\max|\operatorname{Re}\lambda|=0$ 会得到 $13.67$ 并误判"不是 Hopf"，
写成 $\max\operatorname{Re}\lambda=0$ 才对——这是本章第四个陷阱。

\section{蝴蝶效应的定量版}

"蝴蝶效应"通常被讲成一句比喻。本章把它变成一条可以算的曲线：
取两条初值只差 $\varepsilon$ 的轨道，量它们相空间距离 $d(t)$。

\begin{theorem}[可分辨时间]
\label{thm:resolvable}
设 $d(t)=\varepsilon e^{\lambda_1t}$ 直到饱和，显示（或测量）能分辨的最小距离是
$\Delta_{\text{res}}$，则分离第一次被看见的时刻是
\begin{equation}
  t_{\text{vis}}=\frac{1}{\lambda_1}\log\frac{\Delta_{\text{res}}}{\varepsilon}.
  \label{eq:tvis}
\end{equation}
$\varepsilon=10^{-8}$、$\Delta_{\text{res}}=10^{-3}$、$\lambda_1=0.9056$ 给
$t_{\text{vis}}=12.71$。
\end{theorem}

\begin{proof}
解 $\varepsilon e^{\lambda_1t}=\Delta_{\text{res}}$ 得 $t=\lambda_1^{-1}
\log(\Delta_{\text{res}}/\varepsilon)$。$10^{-8}$ 与 $10^{-3}$ 差五个数量级，
$\log10^5=11.5129$，除以 $0.9056$ 得 $12.713$。
\end{proof}

这条公式是本章最有用的一个"防骗"工具：它说明"初值差 $10^{-8}$"这件事
\textbf{在当前显示分辨率下有 $12.7$ 个时间单位的观察窗口}，
而不是"永远看不出来"或者"立刻就乱"。本仓实测（8 对孪生轨道、$h=0.005$）
平均分离第一次越过 $10^{-3}$ 的时刻是 $12.50$，比理论值早 $1.7\%$：
因为平均分离的起点是 $6.0\times10^{-9}$（不是恰好 $10^{-8}$），
且拟合窗口 $[10^{-6},1]$ 内的实测增长率 $0.8876$ 略低于 $\lambda_1$。
两个数都给出，读者自己看差别在哪。

饱和也是可测的：分离不会一直指数增长，它最终停在被吸引子直径决定的量级上。
实测尾段平均分离 $13.87$，单对最大分离 $21.69$。
这条曲线的形状（指数段 + 平台段）就是"可预报时间有限"的全部内容。

\section{周期倍增、Feigenbaum 与"分辨率"}

\begin{theorem}[logistic 的 2 周期与乘子]
\label{thm:two-cycle}
对 $g_r(x)=rx(1-x)$：
\begin{enumerate}
  \item 不动点是 $x=0$ 与 $x=1-1/r$，乘子分别是 $r$ 与 $2-r$；
  \item $r>3$ 时存在 2 周期
        \begin{equation}
          x_\pm=\frac{r+1\pm\sqrt{(r-3)(r+1)}}{2r},
          \label{eq:twocycle}
        \end{equation}
        其乘子是
        \begin{equation}
          M=g_r'(x_+)g_r'(x_-)=r^2(1-2x_+)(1-2x_-)=4+2r-r^2 .
          \label{eq:twomult}
        \end{equation}
\end{enumerate}
\end{theorem}

\begin{proof}
(1) $g_r(x)=x$ 给 $x\bigl(r(1-x)-1\bigr)=0$；$g_r'(x)=r(1-2x)$，代入两个根即得。
(2) 把 $g_r^2(x)-x$ 因式分解。因为 $g_r^2(x)-x$ 一定能被 $g_r(x)-x$ 整除
（不动点也是 2 周期的），余下的是一个二次多项式；直接展开（或长除）得
\[
  \frac{g_r^2(x)-x}{g_r(x)-x}=-r^2x^2+r(r+1)x-(r+1),
\]
它的两个根就是 $x_\pm$（把求根公式写出来即 \eqref{eq:twocycle}，
判别式 $\Delta=r^2(r+1)^2-4r^2(r+1)=r^2(r+1)(r-3)$）。
乘子用链式法则 $M=g_r'(x_+)g_r'(x_-)=r^2\prod_\pm(1-2x_\pm)$；
由韦达定理 $\prod_\pm(1-2x_\pm)=1-2(x_++x_-)+4x_+x_-$，
代入 $x_++x_-=(r+1)/r$、$x_+x_-=(r+1)/r^2$ 得
$M=r^2-2r(r+1)+4(r+1)=4+2r-r^2$。
\end{proof}

\begin{theorem}[倍增点]
\label{thm:doubling}
$1$ 周期在 $r_1=3$ 失稳；2 周期在
\begin{equation}
  r_2=1+\sqrt6=3.449489742783178
  \label{eq:r2}
\end{equation}
失稳（由 $M=-1$）。
\end{theorem}

\begin{proof}
$2-r=-1\Rightarrow r=3$。$4+2r-r^2=-1\Rightarrow r^2-2r-5=0\Rightarrow r=1\pm\sqrt6$，
取 $r>3$ 的根。
\end{proof}

本仓实测：四个倍增点 $3.0000000000000$、$3.4494897427832$、$3.5440903595519$、
$3.5644072660954$ 与闭式/文献值的最大偏差 $1.3\times10^{-15}$；
在每个倍增点上，对应周期的乘子与 $-1$ 的偏差不超过 $1.8\times10^{-14}$。
注意第二条是\textbf{判据本身}的验算：倍增点不是"查表查到的"，
而是"乘子穿过 $-1$"这件事被解出来之后，回代确认的结果。

\begin{theorem}[3 周期窗口的起点是 $1+\sqrt8$]
\label{prop:window3}
周期 3 的窗口从
\begin{equation}
  r_t=1+\sqrt8=3.8284271247461903
  \label{eq:rt}
\end{equation}
开始，那里 $g_r^3(x)-x$ 出现三个二重根
$\approx(0.1599288,\,0.5143551,\,0.9563177)$，且这些点上的乘子等于 $+1$。
$r<r_t$ 时 $g_r^3(x)=x$ 除不动点外没有实根。
\end{theorem}

\begin{proof}[证明与验证]
二重根条件与"乘子 $=+1$"等价：$(g^3)'=g'(x_1)g'(x_2)g'(x_3)=1$。
把 $r=r_t$ 代入 $g_r^3(x)-x$（一个关于 $x$ 的 8 次多项式），
数值求根给出上列三对相差 $<4\times10^{-7}$ 的根
（差值受 8 次多项式在二重根处的条件数限制，不是模型误差），
在这三点上回代 $(g^3)'$ 得 $1\pm4\times10^{-6}$。
本仓另用二分实测窗口起点 $r_t=3.828427125215$，与 $1+\sqrt8$ 差 $4.7\times10^{-10}$。
\end{proof}

窗口有多宽，决定了图要多细才画得出来：实测窗口是
$[3.8284271,\,3.8414990]$，宽 $0.013072$。分岔图的参数网格是
$[2.4,4.0]$ 上 $n_r$ 个点，则窗口里的点数约是
\begin{equation}
  \frac{0.013072}{(4.0-2.4)/(n_r-1)}
  \;\approx\;
  0.00817\,(n_r-1).
  \label{eq:windowres}
\end{equation}
$n_r=700$（\texttt{bifurcation\_}\allowbreak\texttt{logistic} 的默认）给 $5.7$ 个点、
$n_r=600$（\texttt{lyapunov\_}\allowbreak\texttt{overlay} 的默认）给 $4.9$ 个点——
刚好能看出"这里有一小段竖线"；$n_r=200$ 只剩 $1.6$ 个点，
窗口会退化成噪声里的一个毛刺。
\textbf{分辨率不足会把周期窗口画糊}，这是本章第三个陷阱，也是
\texttt{bifurcation\_}\allowbreak\texttt{logistic} 把 $n_r$ 做成参数、
并把"当前窗口里有多少个参数点"写进 \texttt{meta} 的原因。

\begin{theorem}[Feigenbaum 常数是超稳定参数间隔之比的极限]
\label{thm:feigenbaum}
设 $R_n$ 是使 $2^{n-1}$ 周期轨经过临界点 $x=1/2$ 的参数（超稳定参数），则
\begin{equation}
  \delta=\lim_{n\to\infty}\frac{R_n-R_{n-1}}{R_{n+1}-R_n}=4.669201609102990\ldots,
  \qquad
  r_\infty=\lim R_n=3.5699456718\ldots
  \label{eq:delta}
\end{equation}
$\delta$ 没有已知的闭式，它的值只能由序列逼近。
\end{theorem}

\begin{proof}[说明]
超稳定参数由 $g^{2^{n-1}}_{R_n}(1/2)=1/2$ 唯一确定（乘子为 $0$）。
"比值收敛到 $\delta$"这一条属于 Feigenbaum 的重整化理论，本章只用它的结论，
但\textbf{用实测给出收敛速度}：本仓算到 $n=9$ 得
\[
  \delta_n:\ 4.708943,\ 4.680771,\ 4.662960,\ 4.668404,\ 4.668954,\ 4.669157,\ 4.669191,
\]
尾项与 $4.6692016091$ 差 $1.06\times10^{-5}$，而序列中段（第 4 项）的误差是
$7.98\times10^{-4}$——序列确实在收敛，但收敛得慢且不单调。
\textbf{这就是为什么 $\delta$ 的断言只能按实测收敛速度给容差}：
把它写成"等于 $4.6692016091$，容差 $10^{-12}$"是拿序列的一项冒充极限。
\end{proof}

前六个超稳定参数的实测值与文献值的最大偏差 $4.7\times10^{-11}$；
其中 $R_2=1+\sqrt5=3.236067977499790$ 与闭式逐位相符（差 $0$）——
这是本章第五条"闭式对拍"。

\section{Poincare 截面、零倾线与相图}

\begin{definition}[Poincare 截面与返回映射]
取超平面 $\Sigma=\{y:y_{i}=\ell\}$ 与方向 $\pm1$，定义
"沿轨道第一次以规定方向穿过 $\Sigma$"为 $P$。若轨道从 $y\in\Sigma$ 出发、
下一次穿过时的点是 $y'$，则 $P(y)=y'$ 称为返回映射，两次穿越的时间间隔称为返回时间。
\end{definition}

\begin{theorem}[截面上点的定位误差]
\label{thm:poincare}
设步长为 $h$，用"格点两端异号 + 线性插值"定位穿越点时，误差是 $O(h^2)$；
再用截面分量的局部导数 $\dot y_i=f(y)_{i}$ 做一次 Newton 修正
（$t\leftarrow t-y_i/\dot y_i$），误差降到舍入水平。
\end{theorem}

\begin{proof}
设真实穿越时刻是 $t^\ast$，格点下标 $n$ 满足 $t_n\le t^\ast<t_{n+1}$，$h=t_{n+1}-t_n$。
线性插值给出的估计是 $t_n+h\frac{y_i(t_n)}{y_i(t_n)-y_i(t_{n+1})}$。
把 $y_i$ 在 $t^\ast$ 处展开：$y_i(t)=f_i(t^\ast)(t-t^\ast)+O((t-t^\ast)^2)$，
代入得插值估计与 $t^\ast$ 的差为 $O(h^2)$（这是"正割法"的经典结论）。
Newton 一步把线性误差换成二次误差：由 $t^{\text{new}}=t-\frac{y_i(t)}{f_i(t)}$ 与
$y_i(t)=f_i(t^\ast)(t-t^\ast)+O(h^2)$ 得
$t^{\text{new}}-t^\ast=O(h^2)$ 的相对修正，
迭代 $k$ 次后误差是 $O(h^{2^k})$；实现在 $h=0.02$ 上做 $3$ 次即到 $7.3\times10^{-12}$。
\end{proof}

本仓实测（Lorenz，$z=27$ 截面，向上穿越）：
线性插值点到精化点的最大位移按 $h^2$ 下降（$h=0.04/0.02/0.01/0.005$ 给
$0.818/0.184/0.0582/0.0153$），精化后的平面残差 $2.1\times10^{-14}$。
返回时间均值 $0.7487$（$80$ 次穿越）——这是"轨道绕一片翅膀一圈"的时间尺度。

零倾线（nullcline）是 $f$ 的某个分量为零的曲线，两条零倾线的交点就是平衡点。
本章的零倾线用两种独立方式给出：数值追踪（逐列二分 + 把相邻根连成支）
与闭式参数化，两者在 van der Pol 上吻合到 $4.4\times10^{-16}$。
数值追踪有一个必须显式处理的坑：\textbf{连接相邻列的根时用哪个度量}。
van der Pol 的 $\dot y=0$ 零倾线在 $x=\pm1$ 处有竖直渐近线，
$x<-1$ 与 $x>-1$ 上的点纵坐标可以很接近，
只用 $|y|$ 距离连接会把三支接成一支；改用归一化的二维距离
$\sqrt{(\Delta x/w)^2+(\Delta y/h)^2}$ 之后分支数与闭式一致（三支）。

\begin{theorem}[Hopf 标准型的极限环]
\label{thm:hopf-cycle}
对 $\dot x=\mu x-\omega y-x(x^2+y^2)$、$\dot y=\omega x+\mu y-y(x^2+y^2)$：
$\mu>0$ 时原点是不稳定焦点，存在唯一的极限环，它是半径 $\sqrt\mu$ 的圆，
周期精确等于 $2\pi/\omega$。
\end{theorem}

\begin{proof}
转到极坐标 $x=\rho\cos\theta$、$y=\rho\sin\theta$，代入得
\[
  \dot\rho=\mu\rho-\rho^3,\qquad \dot\theta=\omega .
\]
径向方程是一维自治方程，非零平衡点 $\rho=\sqrt\mu$（$\mu>0$）；
$\rho<\sqrt\mu$ 时 $\dot\rho>0$、$\rho>\sqrt\mu$ 时 $\dot\rho<0$，
所以它是全局渐近稳定的（在 $\rho>0$ 上），对应的轨道是半径 $\sqrt\mu$ 的圆。
角速度恒为 $\omega$，故绕一圈的时间是 $2\pi/\omega$。
\end{proof}

本仓实测（$\mu=0.5$）：轨迹尾段平均半径 $0.707106781187$ 与 $\sqrt{0.5}$ 差
$2.2\times10^{-16}$；用截面穿越量出的周期与 $2\pi$ 差 $1.2\times10^{-12}$。
注意 $\mu$ 很小时径向收敛率就是 $\mu$ 本身（$\dot\rho\approx\mu\rho$），
$\mu=0.05$ 要跑到 $t\sim20/\mu=400$ 才看清半径——把时间轴拖长再读，
别把"还没收敛"当成"理论不对"。

\begin{theorem}[van der Pol 在 $\mu\to0$ 的周期修正]
\label{thm:vdp}
$\ddot x+\mu(x^2-1)\dot x+x=0$ 的极限环周期满足
\begin{equation}
  T(\mu)=2\pi\Bigl(1+\frac{\mu^2}{16}+O(\mu^4)\Bigr),
  \qquad\text{即}\qquad
  \lim_{\mu\to0}\frac{T(\mu)-2\pi}{\mu^2}=\frac{2\pi}{16}.
  \label{eq:vdp}
\end{equation}
\end{theorem}

\begin{proof}[证明要点]
用两时间尺度（$t$ 与 $\tau=\mu t$）展开 $x=x_0+\mu x_1+\mu^2x_2+\dots$。
$O(1)$ 给简谐振子 $x_0=A(\tau)\cos t$；$O(\mu)$ 的可解性条件
（消除久期项）给出 $\dot A=A/2-A^3/8$，其非零平衡 $A=2$ 是极限环的振幅；
$O(\mu^2)$ 的可解性条件给出频率修正 $\omega=1-\mu^2/16+O(\mu^4)$，
于是 $T=2\pi/\omega=2\pi(1+\mu^2/16)+O(\mu^4)$。
\end{proof}

本仓实测：$\mu=0$ 的周期 $6.283185307704$ 与 $2\pi$ 差 $5.2\times10^{-10}$；
$\mu=0.02$ 时 $(T-2\pi)/\mu^2=0.392692$，与 $2\pi/16=0.392699081$ 差
$7\times10^{-6}$（相对 $1.8\times10^{-5}$）；$\mu=0.5,1,2$ 时同一比值
单调降到 $0.389962,\,0.380102,\,0.336672$——它确实在往 $\mu\to0$ 的极限走，
而不是"对谁都成立"。

\section{五个可视化该怎么看}

每个帧都把自己的口径写在 \texttt{meta} 与 \texttt{strings} 里。
下面把"该怎么看"和"它承诺了什么"一起列出来。

\begin{enumerate}
  \item \texttt{lorenz\_}\allowbreak\texttt{butterfly}（招牌）：三维吸引子上的两条轨道，
        初值相差 $10^{-8}$。看三件事：两条线什么时候肉眼可分辨（实测 $t_{\text{vis}}=12.50$，
        理论 $12.71$）、副轴的对数分离曲线在哪个窗口里是直线（$[10^{-6},1]$，
        拟合斜率 $0.8876\approx\lambda_1$）、以及它在 $d\approx13.9$ 处变平（饱和）。
        帧里画的是\textbf{一对}孪生轨道（中心轨道 + 它的 $\varepsilon$ 扰动副本，
        两者在 $t=0$ 相距 $10^{-8}$），而副轴上的分离曲线是
        $n_{\text{pairs}}=8$ 对的平均（逐对的曲线也在帧里，
        用来看"一对的运气"与"平均的行为"差多少）；
        另有一条把步长加倍（$h=0.01$）的对照曲线。丢弃暂态 $t=20$，
        样本间隔 $0.05$（$601$ 个采样点）。
        帧里同时给原始 $x/y/z$ 与投影后的 $u/v$，\texttt{strings.projection}
        写着方位角、仰角的定义式——拖动角度参数时前端重画的是同一批三维点，
        \textbf{不是}另算一遍。
  \item \texttt{bifurcation\_}\allowbreak\texttt{logistic}：参数在 $[2.4,4.0]$ 上扫，
        每个参数丢弃 $800$ 次暂态（\texttt{n\_transient}，可调 $100\sim4000$）、
        保留 $64$ 个迭代点（\texttt{n\_keep}）——这就是图上"点云"的密度；
        参数点数 \texttt{n\_r} 与步长 \texttt{dr} 都在 \texttt{meta} 里。时间轴是参数 $r$。
        要看的是：$r=3$ 处一根线裂成两根、$r\approx3.449$ 再裂、
        以及 $r\approx3.83$ 那条**竖直的窄带**（3 周期窗口）。
        分辨率警告写在 \texttt{strings.resolution\_note} 里：默认 $n_r=600$ 时
        $\Delta r=2.67\times10^{-3}$，3 周期窗口（宽 $0.013072$）里只有 $4.9$ 个参数点；
        把它拖到 $200$ 就只剩 $1.6$ 个点，窗口会被画成噪声。
        副轴给的是同一批参数上的 Lyapunov 指数——这张图与下一张是同一份数据。
  \item \texttt{lyapunov\_}\allowbreak\texttt{overlay}：上半是分岔图，
        下半是 $\lambda(r)$。三个解析对拍点被画成叉号：
        $\lambda(2.5)=-0.693147$、$\lambda(3.2)=-0.916291$、$\lambda(4)=0.693148$
        （闭式分别是 $\log0.5$、$\log0.4$、$\log2$）。
        估计口径写在 \texttt{strings.estimate\_}\allowbreak\texttt{protocol}：
        丢弃 $n_{\text{transient}}$ 次暂态、再迭代 $n_{\text{iter}}$ 次、
        按 $n_{\text{block}}$ 块给标准误（默认实测最大标准误 $6.9\times10^{-3}$）。
        要看的是 $\lambda>0$ 的区间与分岔图上"点云糊掉"的区间**重合**，
        以及周期窗口里 $\lambda$ 掉回负值（窗口里的运动是周期的，不分离）。
  \item \texttt{poincare\_}\allowbreak\texttt{section}：三维轨道 + 平面 $z=27$ + 平面上的
        截面点。穿越点数写在 \texttt{meta} 里（\texttt{n\_cross}$=80$）；
        只收向上穿越（$\dot z>0$），所以截面点不会出现"重影"。
        副轴是返回时间序列（\texttt{mean\_return}$=0.7487$）与平面残差：
        \texttt{max\_residual} 应当贴在 $10^{-12}$ 以下（实测 $7.3\times10^{-12}$）；若某帧的残差突然到 $10^{-3}$ 量级，
        那是精化没做够而不是"混沌"——这正是本章第六个陷阱。
  \item \texttt{phase\_}\allowbreak\texttt{portrait}：二维相图上的向量场（箭头）、
        零倾线（两支）、平衡点（相型颜色）与一组初值轨道。
        要看的是"箭头穿过零倾线时哪个分量变号"：$\dot x=0$ 的曲线上箭头是竖直的，
        $\dot y=0$ 的曲线上箭头是水平的，两者交点才是平衡点。
        van der Pol 的 $\dot y=0$ 有竖直渐近线（$x=\pm1$），
        所以那条零倾线是\textbf{三支}而不是一支；
        \texttt{meta} 里的 \texttt{n\_null\_branches}、\texttt{n\_analytic\_branches}
        与 \texttt{max\_residual\_null}（实测 $4.4\times10^{-16}$）就是这件事的对账。
        拖 $\mu$ 过 $0$：原点从稳定焦点变不稳定焦点，极限环出现。
\end{enumerate}

\section{本章的算法路线（P2/P3 合规说明）}

所有结论都由自研代码算出，没有调用任何现成的特征值/QR/最小二乘例程：

\begin{itemize}
  \item \textbf{变分方程 + 自写 QR}：Benettin 算法的核心是每一步对切空间做
        修正 Gram--Schmidt 并累加 $\log|R_{ii}|$。实现是自写的 MGS + 一轮重正交化
        （实测缺陷 $4.4\times10^{-16}$），并有两条恒等式断言：
        $A=QR$ 的重建误差 $<10^{-14}$、$|\det A|=\prod|R_{ii}|$（$<10^{-14}$）；
        Henon 的 $\sum\lambda=\log|b|$ 正是靠第二条成立的。
  \item \textbf{平衡点与零倾线}：特征值由自写的三次/二次求根加 Newton 精化给出；
        零倾线逐列二分求根，再用归一化二维距离把根连成支，
        与闭式参数化对拍（$4.4\times10^{-16}$）。
  \item \textbf{周期轨与乘子}：两条独立路径——Newton 解 $g^p(x)-x=0$
        （带精确周期过滤）与网格扫描 + 向量化二分。两条路径的乘子差 $\le3.5\times10^{-15}$。
        网格扫描有一个必须处理的退化：根正好落在等距节点上时
        （$r=3.2$ 的不动点 $x=0.6875$ 恰好是 $13750/20000$），
        严格的 $F_iF_{i+1}<0$ 判据会漏掉这个根；
        实现里把"端点正好是根"也算作一个括号，并对相邻重复根去重。
  \item \textbf{分岔图与 Lyapunov 曲线}：参数维向量化（一次迭代推进所有参数），
        与逐参数标量路径逐点对拍（偏差 $0.0$）；
        有限时间涨落用分块标准误报告，不用"看起来很平"当证据。
  \item \textbf{拟合}：$\lambda_1$ 的拟合窗口与直线拟合都是自写的
        （正规方程 $2\times2$ 直接解，带 $R^2$），没有用任何现成回归例程。
  \item \textbf{阈值结论}：Hopf 阈值、周期窗口起点、混沌起始点都用二分实测，
        闭式只用来对账（$\rho_H$ 与 $470/19$ 差 $9.8\times10^{-12}$）。
\end{itemize}

% 本文件由 tools/build_book.py 从 dist/reports/assertions.json 生成，请勿手改。
\begin{longtable}{@{}p{0.06\textwidth}p{0.32\textwidth}p{0.22\textwidth}p{0.16\textwidth}@{}}
\caption{章节 3.5 数值断言（生成时间 2026-10-19T16:08:50-0400）}\\
\toprule
类别 & 断言 & 理论值 & 实测值 \\
\midrule
\endfirsthead
\toprule
类别 & 断言 & 理论值 & 实测值 \\
\midrule
\endhead
\bottomrule
\endlastfoot

A & 四个周期倍增点与闭式/文献值的最大偏差 & 周期倍增点由乘子穿过 $-$1 定义：r1=\allowbreak{}3(\allowbreak{}不动点失稳)、\allowbreak{}r2=\allowbreak{}1+$\surd$6、\allowbreak{}r3=\allowbreak{}3.5440903595519225、\allowbreak{}r4=\allowbreak{}3.564407266095432(\allowbreak{}后两个是标准文献值) & $1.332\times 10^{-15}$ \checkmark \\
A & r2 与 1+$\surd$6 的偏差 & 周期倍增点由乘子穿过 $-$1 定义：r1=\allowbreak{}3(\allowbreak{}不动点失稳)、\allowbreak{}r2=\allowbreak{}1+$\surd$6、\allowbreak{}r3=\allowbreak{}3.5440903595519225、\allowbreak{}r4=\allowbreak{}3.564407266095432(\allowbreak{}后两个是标准文献值) & 0 \checkmark \\
A & p=1 在 r1 处的乘子(最接近 $-$1 的那条) & 周期倍增点由乘子穿过 $-$1 定义：r1=\allowbreak{}3(\allowbreak{}不动点失稳)、\allowbreak{}r2=\allowbreak{}1+$\surd$6、\allowbreak{}r3=\allowbreak{}3.5440903595519225、\allowbreak{}r4=\allowbreak{}3.564407266095432(\allowbreak{}后两个是标准文献值) & $4.441\times 10^{-16}$ \checkmark \\
A & p=2 在 r2 处的乘子(最接近 $-$1 的那条) & 周期倍增点由乘子穿过 $-$1 定义：r1=\allowbreak{}3(\allowbreak{}不动点失稳)、\allowbreak{}r2=\allowbreak{}1+$\surd$6、\allowbreak{}r3=\allowbreak{}3.5440903595519225、\allowbreak{}r4=\allowbreak{}3.564407266095432(\allowbreak{}后两个是标准文献值) & $6.661\times 10^{-16}$ \checkmark \\
A & p=4 在 r3 处的乘子(最接近 $-$1 的那条) & 周期倍增点由乘子穿过 $-$1 定义：r1=\allowbreak{}3(\allowbreak{}不动点失稳)、\allowbreak{}r2=\allowbreak{}1+$\surd$6、\allowbreak{}r3=\allowbreak{}3.5440903595519225、\allowbreak{}r4=\allowbreak{}3.564407266095432(\allowbreak{}后两个是标准文献值) & $1.754\times 10^{-14}$ \checkmark \\
A & p=8 在 r4 处的乘子(最接近 $-$1 的那条) & 周期倍增点由乘子穿过 $-$1 定义：r1=\allowbreak{}3(\allowbreak{}不动点失稳)、\allowbreak{}r2=\allowbreak{}1+$\surd$6、\allowbreak{}r3=\allowbreak{}3.5440903595519225、\allowbreak{}r4=\allowbreak{}3.564407266095432(\allowbreak{}后两个是标准文献值) & $1.765\times 10^{-14}$ \checkmark \\
A & r=3.2 的 2 周期点 vs 闭式 & logistic 的 2 周期有闭式：x=\allowbreak{}(\allowbreak{}r+1$\pm$$\surd$(\allowbreak{}(\allowbreak{}r$-$3)(\allowbreak{}r+1)))/\allowbreak{}(\allowbreak{}2r)，乘子 M=\allowbreak{}g'(\allowbreak{}x$_1$)g'(\allowbreak{}x$_2$)=\allowbreak{}4+2r$-$r$^2$ & $1.110\times 10^{-16}$ \checkmark \\
A & r=3.2 的乘子 vs 4+2r$-$r$^2$ & logistic 的 2 周期有闭式：x=\allowbreak{}(\allowbreak{}r+1$\pm$$\surd$(\allowbreak{}(\allowbreak{}r$-$3)(\allowbreak{}r+1)))/\allowbreak{}(\allowbreak{}2r)，乘子 M=\allowbreak{}g'(\allowbreak{}x$_1$)g'(\allowbreak{}x$_2$)=\allowbreak{}4+2r$-$r$^2$ & $6.384\times 10^{-16}$ \checkmark \\
A & r=3.2 的两条乘子路径之差(Newton vs 网格二分) & logistic 的 2 周期有闭式：x=\allowbreak{}(\allowbreak{}r+1$\pm$$\surd$(\allowbreak{}(\allowbreak{}r$-$3)(\allowbreak{}r+1)))/\allowbreak{}(\allowbreak{}2r)，乘子 M=\allowbreak{}g'(\allowbreak{}x$_1$)g'(\allowbreak{}x$_2$)=\allowbreak{}4+2r$-$r$^2$ & $3.469\times 10^{-15}$ \checkmark \\
A & r=3.4 的 2 周期点 vs 闭式 & logistic 的 2 周期有闭式：x=\allowbreak{}(\allowbreak{}r+1$\pm$$\surd$(\allowbreak{}(\allowbreak{}r$-$3)(\allowbreak{}r+1)))/\allowbreak{}(\allowbreak{}2r)，乘子 M=\allowbreak{}g'(\allowbreak{}x$_1$)g'(\allowbreak{}x$_2$)=\allowbreak{}4+2r$-$r$^2$ & $2.220\times 10^{-16}$ \checkmark \\
A & r=3.4 的乘子 vs 4+2r$-$r$^2$ & logistic 的 2 周期有闭式：x=\allowbreak{}(\allowbreak{}r+1$\pm$$\surd$(\allowbreak{}(\allowbreak{}r$-$3)(\allowbreak{}r+1)))/\allowbreak{}(\allowbreak{}2r)，乘子 M=\allowbreak{}g'(\allowbreak{}x$_1$)g'(\allowbreak{}x$_2$)=\allowbreak{}4+2r$-$r$^2$ & $1.332\times 10^{-15}$ \checkmark \\
A & r=3.4 的两条乘子路径之差(Newton vs 网格二分) & logistic 的 2 周期有闭式：x=\allowbreak{}(\allowbreak{}r+1$\pm$$\surd$(\allowbreak{}(\allowbreak{}r$-$3)(\allowbreak{}r+1)))/\allowbreak{}(\allowbreak{}2r)，乘子 M=\allowbreak{}g'(\allowbreak{}x$_1$)g'(\allowbreak{}x$_2$)=\allowbreak{}4+2r$-$r$^2$ & $1.554\times 10^{-15}$ \checkmark \\
A & r=3.23607 的 2 周期点 vs 闭式 & logistic 的 2 周期有闭式：x=\allowbreak{}(\allowbreak{}r+1$\pm$$\surd$(\allowbreak{}(\allowbreak{}r$-$3)(\allowbreak{}r+1)))/\allowbreak{}(\allowbreak{}2r)，乘子 M=\allowbreak{}g'(\allowbreak{}x$_1$)g'(\allowbreak{}x$_2$)=\allowbreak{}4+2r$-$r$^2$ & $1.110\times 10^{-16}$ \checkmark \\
A & r=3.23607 的乘子 vs 4+2r$-$r$^2$ & logistic 的 2 周期有闭式：x=\allowbreak{}(\allowbreak{}r+1$\pm$$\surd$(\allowbreak{}(\allowbreak{}r$-$3)(\allowbreak{}r+1)))/\allowbreak{}(\allowbreak{}2r)，乘子 M=\allowbreak{}g'(\allowbreak{}x$_1$)g'(\allowbreak{}x$_2$)=\allowbreak{}4+2r$-$r$^2$ & $1.437\times 10^{-15}$ \checkmark \\
A & r=3.23607 的两条乘子路径之差(Newton vs 网格二分) & logistic 的 2 周期有闭式：x=\allowbreak{}(\allowbreak{}r+1$\pm$$\surd$(\allowbreak{}(\allowbreak{}r$-$3)(\allowbreak{}r+1)))/\allowbreak{}(\allowbreak{}2r)，乘子 M=\allowbreak{}g'(\allowbreak{}x$_1$)g'(\allowbreak{}x$_2$)=\allowbreak{}4+2r$-$r$^2$ & $2.874\times 10^{-15}$ \checkmark \\
A & 3 周期窗口起点 vs 1+$\surd$8 & 3 周期窗口起点是切分岔：g$^3$(\allowbreak{}x)$-$x 出现二重根，闭式 r=\allowbreak{}1+$\surd$8=\allowbreak{}3.8284271247461903 & $4.685\times 10^{-10}$ \checkmark \\
A & 窗口内外的乘子(内 |M|<1、外无 3 周期) & 3 周期窗口起点是切分岔：g$^3$(\allowbreak{}x)$-$x 出现二重根，闭式 r=\allowbreak{}1+$\surd$8=\allowbreak{}3.8284271247461903 & 0.83293 \checkmark \\
C & 前六个超稳定参数与文献值的最大偏差 & Feigenbaum 常数 $\delta$=\allowbreak{}lim (\allowbreak{}R\_n$-$R\_\{n$-$1\})/\allowbreak{}(\allowbreak{}R\_\{n+1\}$-$R\_n)=\allowbreak{}4.6692016091$\dots$，其中 R\_n 是超稳定参数(\allowbreak{}2\textasciicircum{}\{n$-$1\} 周期轨经过 x=\allowbreak{}1/\allowbreak{}2)。序列收敛不快也不单调，所以容差按实测尾项给 & $4.683\times 10^{-11}$ \checkmark \\
C & R$_2$ 与 1+$\surd$5 的偏差 & Feigenbaum 常数 $\delta$=\allowbreak{}lim (\allowbreak{}R\_n$-$R\_\{n$-$1\})/\allowbreak{}(\allowbreak{}R\_\{n+1\}$-$R\_n)=\allowbreak{}4.6692016091$\dots$，其中 R\_n 是超稳定参数(\allowbreak{}2\textasciicircum{}\{n$-$1\} 周期轨经过 x=\allowbreak{}1/\allowbreak{}2)。序列收敛不快也不单调，所以容差按实测尾项给 & 0 \checkmark \\
C & $\delta$ 序列尾项与 4.6692016091 之差 & Feigenbaum 常数 $\delta$=\allowbreak{}lim (\allowbreak{}R\_n$-$R\_\{n$-$1\})/\allowbreak{}(\allowbreak{}R\_\{n+1\}$-$R\_n)=\allowbreak{}4.6692016091$\dots$，其中 R\_n 是超稳定参数(\allowbreak{}2\textasciicircum{}\{n$-$1\} 周期轨经过 x=\allowbreak{}1/\allowbreak{}2)。序列收敛不快也不单调，所以容差按实测尾项给 & $1.061\times 10^{-5}$ \checkmark \\
C & $\delta$ 序列是否在收敛(尾项误差 < 中段误差) & Feigenbaum 常数 $\delta$=\allowbreak{}lim (\allowbreak{}R\_n$-$R\_\{n$-$1\})/\allowbreak{}(\allowbreak{}R\_\{n+1\}$-$R\_n)=\allowbreak{}4.6692016091$\dots$，其中 R\_n 是超稳定参数(\allowbreak{}2\textasciicircum{}\{n$-$1\} 周期轨经过 x=\allowbreak{}1/\allowbreak{}2)。序列收敛不快也不单调，所以容差按实测尾项给 & $1.061\times 10^{-5}$ \checkmark \\
C & 外推的聚点 r$\infty$ 与文献值 3.5699456718 之差 & Feigenbaum 常数 $\delta$=\allowbreak{}lim (\allowbreak{}R\_n$-$R\_\{n$-$1\})/\allowbreak{}(\allowbreak{}R\_\{n+1\}$-$R\_n)=\allowbreak{}4.6692016091$\dots$，其中 R\_n 是超稳定参数(\allowbreak{}2\textasciicircum{}\{n$-$1\} 周期轨经过 x=\allowbreak{}1/\allowbreak{}2)。序列收敛不快也不单调，所以容差按实测尾项给 & $6.763\times 10^{-11}$ \checkmark \\
A & r=3.2 的最小周期 & logistic 在 r=\allowbreak{}3.2 是 2 周期、\allowbreak{}r=\allowbreak{}3.5 是 4 周期、\allowbreak{}r=\allowbreak{}3.83/\allowbreak{}3.84 落在 3 周期窗口里、\allowbreak{}r=\allowbreak{}4 是混沌(\allowbreak{}$\le$32 步内不闭合)。判据是 |f\textasciicircum{}p(\allowbreak{}x)$-$x|$\le$1e-9，不是肉眼 & 2 \checkmark \\
A & r=3.5 的最小周期 & logistic 在 r=\allowbreak{}3.2 是 2 周期、\allowbreak{}r=\allowbreak{}3.5 是 4 周期、\allowbreak{}r=\allowbreak{}3.83/\allowbreak{}3.84 落在 3 周期窗口里、\allowbreak{}r=\allowbreak{}4 是混沌(\allowbreak{}$\le$32 步内不闭合)。判据是 |f\textasciicircum{}p(\allowbreak{}x)$-$x|$\le$1e-9，不是肉眼 & 4 \checkmark \\
A & r=3.83 的最小周期 & logistic 在 r=\allowbreak{}3.2 是 2 周期、\allowbreak{}r=\allowbreak{}3.5 是 4 周期、\allowbreak{}r=\allowbreak{}3.83/\allowbreak{}3.84 落在 3 周期窗口里、\allowbreak{}r=\allowbreak{}4 是混沌(\allowbreak{}$\le$32 步内不闭合)。判据是 |f\textasciicircum{}p(\allowbreak{}x)$-$x|$\le$1e-9，不是肉眼 & 3 \checkmark \\
A & r=3.84 的最小周期 & logistic 在 r=\allowbreak{}3.2 是 2 周期、\allowbreak{}r=\allowbreak{}3.5 是 4 周期、\allowbreak{}r=\allowbreak{}3.83/\allowbreak{}3.84 落在 3 周期窗口里、\allowbreak{}r=\allowbreak{}4 是混沌(\allowbreak{}$\le$32 步内不闭合)。判据是 |f\textasciicircum{}p(\allowbreak{}x)$-$x|$\le$1e-9，不是肉眼 & 3 \checkmark \\
A & r=4 的周期(应为 0 = 32 步内不闭合) & logistic 在 r=\allowbreak{}3.2 是 2 周期、\allowbreak{}r=\allowbreak{}3.5 是 4 周期、\allowbreak{}r=\allowbreak{}3.83/\allowbreak{}3.84 落在 3 周期窗口里、\allowbreak{}r=\allowbreak{}4 是混沌(\allowbreak{}$\le$32 步内不闭合)。判据是 |f\textasciicircum{}p(\allowbreak{}x)$-$x|$\le$1e-9，不是肉眼 & 0 \checkmark \\
A & 周期窗口里的最大闭合残差 & 周期窗口里闭合残差是舍入水平(\allowbreak{}1e-15)，混沌区是 O(\allowbreak{}1e-2)：两类点必须分在对数轴两端；混沌区的残差**不能**是 0(\allowbreak{}那会把混沌画成完美闭合) & $9.279\times 10^{-10}$ \checkmark \\
A & 混沌区闭合残差的中位数 & 周期窗口里闭合残差是舍入水平(\allowbreak{}1e-15)，混沌区是 O(\allowbreak{}1e-2)：两类点必须分在对数轴两端；混沌区的残差**不能**是 0(\allowbreak{}那会把混沌画成完美闭合) & 0.005202 \checkmark \\
A & 混沌区里残差 > 1e-4 的参数占比 & 周期窗口里闭合残差是舍入水平(\allowbreak{}1e-15)，混沌区是 O(\allowbreak{}1e-2)：两类点必须分在对数轴两端；混沌区的残差**不能**是 0(\allowbreak{}那会把混沌画成完美闭合) & 0.943503 \checkmark \\
A & 向量化 vs 标量的周期差异计数 & 向量化的 detect\_periods 与逐参数标量 detect\_period 必须给出同一个周期与同一个残差(\allowbreak{}批量路径只是把参数维摊开；残差允许舍入量级的差) & 0 \checkmark \\
A & 向量化 vs 标量的残差最大差 & 向量化的 detect\_periods 与逐参数标量 detect\_period 必须给出同一个周期与同一个残差(\allowbreak{}批量路径只是把参数维摊开；残差允许舍入量级的差) & $3.220\times 10^{-15}$ \checkmark \\
A & r=2.5 的不动点是否稳定 & r=\allowbreak{}4 的超稳定不动点(\allowbreak{}x=\allowbreak{}1/\allowbreak{}2 是 r=\allowbreak{}2 的临界点)与周期窗口的存在性给出 Logistic 映射分岔结构的两个锚点：r<3 只有 1 周期、\allowbreak{}r=\allowbreak{}3.2 的乘子为 0.16 & 1 \checkmark \\
A & r=3.2 的不动点是否稳定 & r=\allowbreak{}4 的超稳定不动点(\allowbreak{}x=\allowbreak{}1/\allowbreak{}2 是 r=\allowbreak{}2 的临界点)与周期窗口的存在性给出 Logistic 映射分岔结构的两个锚点：r<3 只有 1 周期、\allowbreak{}r=\allowbreak{}3.2 的乘子为 0.16 & 0 \checkmark \\
A & r=3.2 处不动点的乘子里有没有 2$-$r=$-$1.2 & r=\allowbreak{}4 的超稳定不动点(\allowbreak{}x=\allowbreak{}1/\allowbreak{}2 是 r=\allowbreak{}2 的临界点)与周期窗口的存在性给出 Logistic 映射分岔结构的两个锚点：r<3 只有 1 周期、\allowbreak{}r=\allowbreak{}3.2 的乘子为 0.16 & 0 \checkmark \\
A & r=3.2 处稳定 2 周期的乘子 vs 4+2r$-$r$^2$=0.16 & r=\allowbreak{}4 的超稳定不动点(\allowbreak{}x=\allowbreak{}1/\allowbreak{}2 是 r=\allowbreak{}2 的临界点)与周期窗口的存在性给出 Logistic 映射分岔结构的两个锚点：r<3 只有 1 周期、\allowbreak{}r=\allowbreak{}3.2 的乘子为 0.16 & $2.276\times 10^{-15}$ \checkmark \\
A & 三个平衡点的 |f| 残差 & Lorenz 的非平凡平衡点 C$\pm$=\allowbreak{}(\allowbreak{}$\pm$$\surd$(\allowbreak{}$\beta$(\allowbreak{}$\rho$$-$1)),\allowbreak{} $\pm$$\surd$(\allowbreak{}$\beta$(\allowbreak{}$\rho$$-$1)),\allowbreak{} $\rho$$-$1) 与原点都是 f=\allowbreak{}0 的精确解；散度 $\nabla$$\cdot$f=\allowbreak{}$-$(\allowbreak{}$\sigma$+1+$\beta$) 与状态无关(\allowbreak{}常数) & $1.421\times 10^{-14}$ \checkmark \\
A & C$\pm$ 与闭式 $\surd$($\beta$($\rho$$-$1)) 的偏差 & Lorenz 的非平凡平衡点 C$\pm$=\allowbreak{}(\allowbreak{}$\pm$$\surd$(\allowbreak{}$\beta$(\allowbreak{}$\rho$$-$1)),\allowbreak{} $\pm$$\surd$(\allowbreak{}$\beta$(\allowbreak{}$\rho$$-$1)),\allowbreak{} $\rho$$-$1) 与原点都是 f=\allowbreak{}0 的精确解；散度 $\nabla$$\cdot$f=\allowbreak{}$-$(\allowbreak{}$\sigma$+1+$\beta$) 与状态无关(\allowbreak{}常数) & 0 \checkmark \\
A & 散度 = $-$($\sigma$+1+$\beta$) & Lorenz 的非平凡平衡点 C$\pm$=\allowbreak{}(\allowbreak{}$\pm$$\surd$(\allowbreak{}$\beta$(\allowbreak{}$\rho$$-$1)),\allowbreak{} $\pm$$\surd$(\allowbreak{}$\beta$(\allowbreak{}$\rho$$-$1)),\allowbreak{} $\rho$$-$1) 与原点都是 f=\allowbreak{}0 的精确解；散度 $\nabla$$\cdot$f=\allowbreak{}$-$(\allowbreak{}$\sigma$+1+$\beta$) 与状态无关(\allowbreak{}常数) & 0 \checkmark \\
A & 原点的三个实特征值 & Lorenz 的非平凡平衡点 C$\pm$=\allowbreak{}(\allowbreak{}$\pm$$\surd$(\allowbreak{}$\beta$(\allowbreak{}$\rho$$-$1)),\allowbreak{} $\pm$$\surd$(\allowbreak{}$\beta$(\allowbreak{}$\rho$$-$1)),\allowbreak{} $\rho$$-$1) 与原点都是 f=\allowbreak{}0 的精确解；散度 $\nabla$$\cdot$f=\allowbreak{}$-$(\allowbreak{}$\sigma$+1+$\beta$) 与状态无关(\allowbreak{}常数) & $5.488\times 10^{-7}$ \checkmark \\
A & C$\pm$ 复共轭对的虚部与 10.194501 之差 & Lorenz 的非平凡平衡点 C$\pm$=\allowbreak{}(\allowbreak{}$\pm$$\surd$(\allowbreak{}$\beta$(\allowbreak{}$\rho$$-$1)),\allowbreak{} $\pm$$\surd$(\allowbreak{}$\beta$(\allowbreak{}$\rho$$-$1)),\allowbreak{} $\rho$$-$1) 与原点都是 f=\allowbreak{}0 的精确解；散度 $\nabla$$\cdot$f=\allowbreak{}$-$(\allowbreak{}$\sigma$+1+$\beta$) 与状态无关(\allowbreak{}常数) & $4.221\times 10^{-6}$ \checkmark \\
A & $\rho$=28 处 C$\pm$ 的相型(鞍焦) & Lorenz 的非平凡平衡点 C$\pm$=\allowbreak{}(\allowbreak{}$\pm$$\surd$(\allowbreak{}$\beta$(\allowbreak{}$\rho$$-$1)),\allowbreak{} $\pm$$\surd$(\allowbreak{}$\beta$(\allowbreak{}$\rho$$-$1)),\allowbreak{} $\rho$$-$1) 与原点都是 f=\allowbreak{}0 的精确解；散度 $\nabla$$\cdot$f=\allowbreak{}$-$(\allowbreak{}$\sigma$+1+$\beta$) 与状态无关(\allowbreak{}常数) & 7 \checkmark \\
C & 二分实测 $\rho$\_\allowbreak{}H 与 470/19 & Hopf 分岔的定义是平衡点的一对特征值穿过虚轴：max Re $\lambda$(\allowbreak{}$\rho$) 在 $\rho$\_H 变号。二分实测该变号点，与闭式 $\rho$\_H=\allowbreak{}$\sigma$(\allowbreak{}$\sigma$+$\beta$+3)/\allowbreak{}(\allowbreak{}$\sigma$$-$$\beta$$-$1)=\allowbreak{}470/\allowbreak{}19 对拍 & $9.809\times 10^{-12}$ \checkmark \\
C & $\rho$\_\allowbreak{}H 处 C$\pm$ 的 max Re $\lambda$(应为 0) & Hopf 分岔的定义是平衡点的一对特征值穿过虚轴：max Re $\lambda$(\allowbreak{}$\rho$) 在 $\rho$\_H 变号。二分实测该变号点，与闭式 $\rho$\_H=\allowbreak{}$\sigma$(\allowbreak{}$\sigma$+$\beta$+3)/\allowbreak{}(\allowbreak{}$\sigma$$-$$\beta$$-$1)=\allowbreak{}470/\allowbreak{}19 对拍 & $5.928\times 10^{-16}$ \checkmark \\
C & $\rho$\_\allowbreak{}H 处第三个特征值(负实根，与 $-$($\sigma$+1+$\beta$) 同值) & Hopf 分岔的定义是平衡点的一对特征值穿过虚轴：max Re $\lambda$(\allowbreak{}$\rho$) 在 $\rho$\_H 变号。二分实测该变号点，与闭式 $\rho$\_H=\allowbreak{}$\sigma$(\allowbreak{}$\sigma$+$\beta$+3)/\allowbreak{}(\allowbreak{}$\sigma$$-$$\beta$$-$1)=\allowbreak{}470/\allowbreak{}19 对拍 & $1.776\times 10^{-15}$ \checkmark \\
A & van der Pol 的 dy/dt=0 零倾线分支数(数值 vs 闭式) & van der Pol 的 dy/\allowbreak{}dt=\allowbreak{}0 零倾线有闭式 y=\allowbreak{}x/\allowbreak{}(\allowbreak{}$\mu$(\allowbreak{}1$-$x$^2$))；数值追踪(\allowbreak{}逐列二分 + 归一化二维距离连接)必须给出同一条曲线，且分支数一致 & 3 \checkmark \\
A & 数值零倾线 vs 闭式 y=x/($\mu$(1$-$x$^2$)) 的最大偏差 & van der Pol 的 dy/\allowbreak{}dt=\allowbreak{}0 零倾线有闭式 y=\allowbreak{}x/\allowbreak{}(\allowbreak{}$\mu$(\allowbreak{}1$-$x$^2$))；数值追踪(\allowbreak{}逐列二分 + 归一化二维距离连接)必须给出同一条曲线，且分支数一致 & $4.441\times 10^{-16}$ \checkmark \\
A & 跨过 x=$-$1 渐近线的分支数(应为 0) & van der Pol 的 dy/\allowbreak{}dt=\allowbreak{}0 零倾线有闭式 y=\allowbreak{}x/\allowbreak{}(\allowbreak{}$\mu$(\allowbreak{}1$-$x$^2$))；数值追踪(\allowbreak{}逐列二分 + 归一化二维距离连接)必须给出同一条曲线，且分支数一致 & 0 \checkmark \\
A & 阻尼摆 dy/dt=0 零倾线 vs 闭式 y=$-$sin(x)/c & van der Pol 的 dy/\allowbreak{}dt=\allowbreak{}0 零倾线有闭式 y=\allowbreak{}x/\allowbreak{}(\allowbreak{}$\mu$(\allowbreak{}1$-$x$^2$))；数值追踪(\allowbreak{}逐列二分 + 归一化二维距离连接)必须给出同一条曲线，且分支数一致 & $4.441\times 10^{-16}$ \checkmark \\
A & 截面点离平面的最大距离(精化后) & 定向 Poincare 截面只收单向穿越：截面上每一点都必须满足 f(\allowbreak{}y)[平面分量]>0(\allowbreak{}向上穿越)且落在平面上；返回时间与穿越顺序一一对应 & $7.280\times 10^{-12}$ \checkmark \\
A & 穿越点上的 dz/dt 全为正(定向判据) & 定向 Poincare 截面只收单向穿越：截面上每一点都必须满足 f(\allowbreak{}y)[平面分量]>0(\allowbreak{}向上穿越)且落在平面上；返回时间与穿越顺序一一对应 & 69.0878 \checkmark \\
A & 相邻穿越的时间间隔(均值$\pm$极差) & 定向 Poincare 截面只收单向穿越：截面上每一点都必须满足 f(\allowbreak{}y)[平面分量]>0(\allowbreak{}向上穿越)且落在平面上；返回时间与穿越顺序一一对应 & 0.748659 \checkmark \\
C & 精化位移随步长的收敛(最大位移) & 截面点的定位误差有两个来源：线性插值是 O(\allowbreak{}h$^2$)，Newton 精化把它压到舍入水平。「线性插值点 $\to$ 精化点」的位移必须按 h$^2$ 收敛 & 0.0153331 \checkmark \\
C & 线性插值定位误差的收敛阶 & 2 & 1.88789 \checkmark \\
C & 平面残差(Newton 精化后) & 截面点的定位误差有两个来源：线性插值是 O(\allowbreak{}h$^2$)，Newton 精化把它压到舍入水平。「线性插值点 $\to$ 精化点」的位移必须按 h$^2$ 收敛 & $1.421\times 10^{-14}$ \checkmark \\
A & $\mu$=0 的周期 $-$ 2$\pi$ & $\mu$=\allowbreak{}0 的 van der Pol 退化成简谐振子 x''+x=\allowbreak{}0，周期精确等于 2$\pi$；$\mu$$\to$0 时极限环周期的修正是 (\allowbreak{}T$-$2$\pi$)/\allowbreak{}$\mu$$^2$ $\to$ 2$\pi$/\allowbreak{}16 & $5.243\times 10^{-10}$ \checkmark \\
A & (T$-$2$\pi$)/$\mu$$^2$ 在 $\mu$=0.02 处的值 vs 2$\pi$/16 & $\mu$=\allowbreak{}0 的 van der Pol 退化成简谐振子 x''+x=\allowbreak{}0，周期精确等于 2$\pi$；$\mu$$\to$0 时极限环周期的修正是 (\allowbreak{}T$-$2$\pi$)/\allowbreak{}$\mu$$^2$ $\to$ 2$\pi$/\allowbreak{}16 & $6.716\times 10^{-6}$ \checkmark \\
A & 极限环振幅随 $\mu$ 的变化(2 的修正) & $\mu$=\allowbreak{}0 的 van der Pol 退化成简谐振子 x''+x=\allowbreak{}0，周期精确等于 2$\pi$；$\mu$$\to$0 时极限环周期的修正是 (\allowbreak{}T$-$2$\pi$)/\allowbreak{}$\mu$$^2$ $\to$ 2$\pi$/\allowbreak{}16 & 0.00861986 \checkmark \\
A & 极限环半径 vs $\surd$$\mu$ & Hopf 标准型 $\rho$'=\allowbreak{}$\mu$$\rho$$-$$\rho$$^3$、\allowbreak{}$\theta$'=\allowbreak{}$\omega$ 的极限环是半径 $\surd$$\mu$ 的圆，周期精确等于 2$\pi$/\allowbreak{}$\omega$ & $2.220\times 10^{-16}$ \checkmark \\
A & 极限环周期 vs 2$\pi$/$\omega$ & Hopf 标准型 $\rho$'=\allowbreak{}$\mu$$\rho$$-$$\rho$$^3$、\allowbreak{}$\theta$'=\allowbreak{}$\omega$ 的极限环是半径 $\surd$$\mu$ 的圆，周期精确等于 2$\pi$/\allowbreak{}$\omega$ & $8.517\times 10^{-12}$ \checkmark \\
A & 批量积分 vs 单条积分(逐点最大差) & integrate\_ensemble 把 (\allowbreak{}n,\allowbreak{}dim) 摊平后用同一个 rk\_step 推进；n=\allowbreak{}1 时它必须与单条轨迹的 integrate\_fixed 逐点相同(\allowbreak{}两条独立的代码路径) & 0 \checkmark \\
A & lorenz\_\allowbreak{}butterfly：由 x/y/z + strings 角度重算 u/v 的偏差 & 三维帧必须同时给原始 x/\allowbreak{}y/\allowbreak{}z 与投影后的 u/\allowbreak{}v，并把投影参数写进 strings；帧里的 u/\allowbreak{}v 必须能由同帧的 x/\allowbreak{}y/\allowbreak{}z 与同一组角度重算出来 & 0 \checkmark \\
A & poincare\_\allowbreak{}section：由 x/y/z + strings 角度重算 u/v 的偏差 & 三维帧必须同时给原始 x/\allowbreak{}y/\allowbreak{}z 与投影后的 u/\allowbreak{}v，并把投影参数写进 strings；帧里的 u/\allowbreak{}v 必须能由同帧的 x/\allowbreak{}y/\allowbreak{}z 与同一组角度重算出来 & 0 \checkmark \\
A & meta.t\_\allowbreak{}visible vs 数组里第一次越过 1e-3 的时刻 & 帧里报出来的标量必须与同帧的数组一致(\allowbreak{}t\_visible 是 sep\_mean 第一次达到显示分辨率１e-3 的时刻)------否则前端画的与 meta 说的是两回事 & 0 \checkmark \\
A & 孪生实验的拟合增长率(默认参数) & 帧里报出来的标量必须与同帧的数组一致(\allowbreak{}t\_visible 是 sep\_mean 第一次达到显示分辨率１e-3 的时刻)------否则前端画的与 meta 说的是两回事 & 0.887642 \checkmark \\
A & h 与 2h 两条分离曲线的增长率之差(默认参数) & 帧里报出来的标量必须与同帧的数组一致(\allowbreak{}t\_visible 是 sep\_mean 第一次达到显示分辨率１e-3 的时刻)------否则前端画的与 meta 说的是两回事 & 0.00195081 \checkmark \\
A & 第 1 条初值的 $\Sigma$$\lambda$ 与散度之差 & Lyapunov 指数的和等于散度的时间平均：$\Sigma$$\lambda$\_i =\allowbreak{} $\langle$$\nabla$$\cdot$f$\rangle$。Lorenz 的散度是常数 $-$(\allowbreak{}$\sigma$+1+$\beta$)=\allowbreak{}$-$13.6667，所以这条恒等式给出一个不依赖任何文献值的标尺 & 0.000102712 \checkmark \\
A & 第 1 条初值的 $\lambda$$_2$(流向，应趋于 0) & Lyapunov 指数的和等于散度的时间平均：$\Sigma$$\lambda$\_i =\allowbreak{} $\langle$$\nabla$$\cdot$f$\rangle$。Lorenz 的散度是常数 $-$(\allowbreak{}$\sigma$+1+$\beta$)=\allowbreak{}$-$13.6667，所以这条恒等式给出一个不依赖任何文献值的标尺 & 0.000324632 \checkmark \\
A & 第 2 条初值的 $\Sigma$$\lambda$ 与散度之差 & Lyapunov 指数的和等于散度的时间平均：$\Sigma$$\lambda$\_i =\allowbreak{} $\langle$$\nabla$$\cdot$f$\rangle$。Lorenz 的散度是常数 $-$(\allowbreak{}$\sigma$+1+$\beta$)=\allowbreak{}$-$13.6667，所以这条恒等式给出一个不依赖任何文献值的标尺 & 0.000101611 \checkmark \\
A & 第 2 条初值的 $\lambda$$_2$(流向，应趋于 0) & Lyapunov 指数的和等于散度的时间平均：$\Sigma$$\lambda$\_i =\allowbreak{} $\langle$$\nabla$$\cdot$f$\rangle$。Lorenz 的散度是常数 $-$(\allowbreak{}$\sigma$+1+$\beta$)=\allowbreak{}$-$13.6667，所以这条恒等式给出一个不依赖任何文献值的标尺 & 0.00062247 \checkmark \\
A & Gram--Schmidt 的最大正交性缺陷 & Lyapunov 指数的和等于散度的时间平均：$\Sigma$$\lambda$\_i =\allowbreak{} $\langle$$\nabla$$\cdot$f$\rangle$。Lorenz 的散度是常数 $-$(\allowbreak{}$\sigma$+1+$\beta$)=\allowbreak{}$-$13.6667，所以这条恒等式给出一个不依赖任何文献值的标尺 & $6.661\times 10^{-16}$ \checkmark \\
A & Henon 的 $\Sigma$$\lambda$ 与 log|b| 之差 & Henon 映射的 Jacobi 行列式处处等于 $-$b，所以两个 Lyapunov 指数之和精确等于 log|b|(\allowbreak{}b=\allowbreak{}0.3 $\to$ $-$1.2039728043259361)，与轨道无关 & $1.177\times 10^{-14}$ \checkmark \\
A & $\lambda$$_1$ 与文献值 0.41922 之差 & Henon 映射的 Jacobi 行列式处处等于 $-$b，所以两个 Lyapunov 指数之和精确等于 log|b|(\allowbreak{}b=\allowbreak{}0.3 $\to$ $-$1.2039728043259361)，与轨道无关 & 0.00038919 \checkmark \\
A & $\lambda$(2.5) 与闭式之差 & logistic 映射的 Lyapunov 指数有三处闭式：1<r<3 时 $\lambda$=\allowbreak{}log|2$-$r|；2 周期窗口内 $\lambda$=\allowbreak{}$1/\allowbreak{}2$log|4+2r$-$r$^2$|；r=\allowbreak{}4 共轭到 Bernoulli 移位，$\lambda$=\allowbreak{}log2 & $3.819\times 10^{-14}$ \checkmark \\
A & $\lambda$(3.2) 与闭式之差 & logistic 映射的 Lyapunov 指数有三处闭式：1<r<3 时 $\lambda$=\allowbreak{}log|2$-$r|；2 周期窗口内 $\lambda$=\allowbreak{}$1/\allowbreak{}2$log|4+2r$-$r$^2$|；r=\allowbreak{}4 共轭到 Bernoulli 移位，$\lambda$=\allowbreak{}log2 & $1.088\times 10^{-14}$ \checkmark \\
A & $\lambda$(4.0) 与闭式之差 & logistic 映射的 Lyapunov 指数有三处闭式：1<r<3 时 $\lambda$=\allowbreak{}log|2$-$r|；2 周期窗口内 $\lambda$=\allowbreak{}$1/\allowbreak{}2$log|4+2r$-$r$^2$|；r=\allowbreak{}4 共轭到 Bernoulli 移位，$\lambda$=\allowbreak{}log2 & $1.847\times 10^{-6}$ \checkmark \\
A & $\lambda$$_1$ 的均值与公认值 0.9056 之差 & Lorenz(\allowbreak{}$\sigma$=\allowbreak{}10,\allowbreak{}$\rho$=\allowbreak{}28,\allowbreak{}$\beta$=\allowbreak{}8/\allowbreak{}3) 的最大 Lyapunov 指数公认值 0.9056(\allowbreak{}Sprott 2003 表 5.1)。有限时间估计 T=\allowbreak{}2000 的涨落由跨初值的标准差实测，容差据此给出 & 0.00154537 \checkmark \\
A & 公认值落在实测的 $\pm$2$\sigma$ 区间内 & Lorenz(\allowbreak{}$\sigma$=\allowbreak{}10,\allowbreak{}$\rho$=\allowbreak{}28,\allowbreak{}$\beta$=\allowbreak{}8/\allowbreak{}3) 的最大 Lyapunov 指数公认值 0.9056(\allowbreak{}Sprott 2003 表 5.1)。有限时间估计 T=\allowbreak{}2000 的涨落由跨初值的标准差实测，容差据此给出 & 0.150413 \checkmark \\
C & h 减半后 $\lambda$$_1$ 的变化 & 定步长积分的 Lyapunov 估计必须随步长收敛：h 减半后 $\lambda$$_1$ 的变化量要小于积分误差允许的量级，否则「混沌」可能只是离散化噪声 & 0.0143774 \checkmark \\
C & h=0.01 的 $\Sigma$$\lambda$ 与散度之差 & 定步长积分的 Lyapunov 估计必须随步长收敛：h 减半后 $\lambda$$_1$ 的变化量要小于积分误差允许的量级，否则「混沌」可能只是离散化噪声 & 0.000103136 \checkmark \\
C & h=0.005 的 $\Sigma$$\lambda$ 与散度之差 & 定步长积分的 Lyapunov 估计必须随步长收敛：h 减半后 $\lambda$$_1$ 的变化量要小于积分误差允许的量级，否则「混沌」可能只是离散化噪声 & $6.215\times 10^{-6}$ \checkmark \\
C & 孪生分离曲线的增长率(h 与 2h)之差 & 定步长积分的 Lyapunov 估计必须随步长收敛：h 减半后 $\lambda$$_1$ 的变化量要小于积分误差允许的量级，否则「混沌」可能只是离散化噪声 & 0.00277876 \checkmark \\
C & QR 间隔 10 vs 20 的谱差 & QR 重正交化的间隔只影响浮点误差与代价，不影响数学：间隔 10 与 20 必须给出同一个谱 & $2.132\times 10^{-14}$ \checkmark \\
C & 孪生增长率 vs Benettin $\lambda$$_1$ & 孪生轨迹的有限时间增长率必须与 Benettin 谱给出的 $\lambda$$_1$ 同量级(\allowbreak{}两者是同一个量的两种测法：一个直接量两条轨道之差，一个走变分方程 + QR) & 0.0236775 \checkmark \\
C & 饱和距离(尾段平均分离) & 孪生轨迹的有限时间增长率必须与 Benettin 谱给出的 $\lambda$$_1$ 同量级(\allowbreak{}两者是同一个量的两种测法：一个直接量两条轨道之差，一个走变分方程 + QR) & 13.8653 \checkmark \\
C & $\lambda$(4) 与 ln2 之差 vs 分块标准误 & 有限时间 $\lambda$ 估计的误差应当落在分块标准误的量级内：块标准误度量的是同一件事(\allowbreak{}有限时间涨落)，如果误差远大于它，说明估计有系统偏差 & 0.219304 \checkmark \\
C & Kaplan--Yorke 维数(Lorenz，由实测谱推出) & 有限时间 $\lambda$ 估计的误差应当落在分块标准误的量级内：块标准误度量的是同一件事(\allowbreak{}有限时间涨落)，如果误差远大于它，说明估计有系统偏差 & 2.06249 \checkmark \\
A & A=QR 的重建误差 & 自写的修正 Gram--Schmidt 必须满足 A=\allowbreak{}QR 与 Q\textasciicircum{}T Q=\allowbreak{}I；det 的绝对值等于 $\Pi$|R\_ii| & $1.110\times 10^{-16}$ \checkmark \\
A & |det A| vs $\Pi$|R\_\allowbreak{}ii| & 自写的修正 Gram--Schmidt 必须满足 A=\allowbreak{}QR 与 Q\textasciicircum{}T Q=\allowbreak{}I；det 的绝对值等于 $\Pi$|R\_ii| & $8.882\times 10^{-16}$ \checkmark \\
A & 数值散度 vs 闭式 $-$13.6667 & 自写的修正 Gram--Schmidt 必须满足 A=\allowbreak{}QR 与 Q\textasciicircum{}T Q=\allowbreak{}I；det 的绝对值等于 $\Pi$|R\_ii| & $1.466\times 10^{-9}$ \checkmark \\
C & 向量化分岔图 vs 标量迭代的最大偏差 & 分岔图的点云与逐参数标量迭代必须给出同一批点(\allowbreak{}向量化路径不能改变结果) & 0 \checkmark \\
\end{longtable}


\section{它会怎样骗你（本章版）}

\begin{enumerate}
  \item \textbf{把"画得像蝴蝶"当成混沌的证据}：吸引子的形状与"对初值敏感"
        是两件事。本章的判据是"孪生分离的拟合增长率 + 步长减半不变"，
        图上像不像只负责好看。
  \item \textbf{不给步长证据就宣布混沌}：$h$ 太大时截断误差自己就会指数放大。
        本章的口径是：$h$ 从 $0.01$ 减半到 $0.005$，Benettin 的 $\lambda_1$ 从
        $0.917925$ 变到 $0.903548$（差 $0.0144$，$T=400$），
        孪生曲线增长率从 $0.8908$ 变到 $0.8881$（差 $0.0028$，$8$ 对平均），
        $\sum\lambda$ 与散度的差在两种步长下分别是 $1.0\times10^{-4}$ 与
        $6.2\times10^{-6}$。
        没有这三条，"混沌"可能只是离散化噪声。
  \item \textbf{分辨率不足会把周期窗口画糊}：3 周期窗口宽 $0.013072$。
        $n_r=600$ 时窗口里 $4.9$ 个参数点，$n_r=200$ 时 $1.6$ 个。
        图上"这里好像有一竖条"很可能只是采样点太少；
        本章把"窗口里的参数点数"直接写进 \texttt{meta}。
  \item \textbf{把 $\max|\operatorname{Re}\lambda|=0$ 当 Hopf 判据}：
        $\rho_H$ 处 $C_\pm$ 的第三个特征值是 $-13.6667$ 而不是 $0$，
        写成绝对值会得到 $13.67$ 并得出"没有 Hopf"的相反结论。判据是
        $\max\operatorname{Re}\lambda$（带符号）。
  \item \textbf{把鞍焦当鞍点}：$\rho=28$ 的 $C_\pm$ 是
        $(-13.8546,\,0.0940\pm10.1945\mathrm{i})$：一个稳定实方向 + 一个不稳定螺旋面。
        本仓初版把它判成 \texttt{saddle}，于是"两片翅膀"在相图里失去了来源。
  \item \textbf{把截面的线性插值点当交点}：$h=0.02$ 时线性插值点的平面残差是
        $3.4\times10^{-3}$，看上去"截面点散成一条带"；
        这既不是混沌也不是数值不稳定，而是定位误差 $O(h^2)$。
        做一次 Newton 修正后残差到 $10^{-12}$，带就收成点。
  \item \textbf{把混沌区的"闭合残差"填成 0}：本章初版对检测不到周期的参数
        直接把残差写成 $0.0$，图上"混沌"看起来完美闭合。
        现在的实现一律给 $\min_{p\le P}|g^p(x)-x|$：窗口里 $<10^{-9}$、
        混沌区中位 $5.2\times10^{-3}$（$94.4\%$ 大于 $10^{-4}$）。
        另有一条"灰区"要记住：固定容差必然造出一条近闭合轨道带
        （实测最小 $2.1\times10^{-9}$），所以 \texttt{meta} 里同时给出用了多大的容差。
  \item \textbf{把 Feigenbaum 常数当"算出来就等于"}：$\delta$ 是序列的极限，
        本仓算到 $n=9$ 的尾项与文献值差 $1.06\times10^{-5}$，中段差 $7.98\times10^{-4}$。
        给 $\delta$ 的断言只能按实测收敛速度给容差；把容差写成 $10^{-12}$
        等于假装序列已经到极限。
  \item \textbf{把有限时间 Lyapunov 估计当极限值}：同一组参数、不同初值、
        $T=2000$ 的 $\lambda_1$ 标准差是 $0.0103$；把 $T$ 从 $400$ 加到 $2000$，
        均值从 $0.9017$ 爬到 $0.9041$。所以"$\lambda_1=0.9056$"这句话
        必须带容差与口径，本章给的是 $\pm0.02$（约 $2\sigma$）。
  \item \textbf{把饱和后的分离当发散}：$d(t)$ 在 $13.9$ 附近变平，
        这不是"数值解崩了"，而是吸引子的直径。往上翻十倍去看
        $\log_{10}d$ 曲线，指数段只占前 $40\%$ 的时间。
  \item \textbf{把时间轴当时间}：\texttt{bifurcation\_}\allowbreak\texttt{logistic}
        与 \texttt{lyapunov\_}\allowbreak\texttt{overlay} 的时间轴是\textbf{参数} $r$，
        \texttt{lorenz\_}\allowbreak\texttt{butterfly} 是真实时间但只画到 $t=30$
        （样本间隔 $0.05$），\texttt{poincare\_}\allowbreak\texttt{section}
        画到 $t=60$。每帧的 \texttt{meta} 都给出当前口径，拖到头是夹住，不是卡住。
  \item \textbf{把投影后的二维图当三维结构}：$u/v$ 是旋转投影，
        "两片翅膀重叠"可能只是视角问题。帧里同时给原始 $x/y/z$，
        拖 \texttt{azimuth}/\texttt{elevation} 看同一批点的另一个视角——
        如果那是同一条轨道，换个角度它仍然是一条轨道。
\end{enumerate}
